\chapter{Regular Conditional Laws and Conditioning Coordinates}
\label{sec:p15-regular-conditional-laws}

A scalar conditional probability is specified separately for each event.
A regular conditional law supplies one probability measure for all target
events at each conditioning argument. The source is Durrett's
Section~4.1.3, with the sphere example in Section~4.1.1 and the conditional
H\"older exercise \cite[pp.~210, 214--216]{DurrettPTEAuthor2019}.
The construction below uses the proved conditional expectations, Stieltjes
probabilities and Borel real codings.

\section{Probability kernels and their integrals}
\label{sec:p15-conditional-kernels}

\begin{hbdefinition}{def:p15-probability-kernel}{A probability kernel}
For measurable spaces \((T,\mathcal T)\) and \((S,\mathcal S)\), a
probability kernel is a function
\(K:T\times\mathcal S\to[0,1]\) such that:
\begin{enumerate}
\item \(y\mapsto K(y,A)\) is \(\mathcal T\)-measurable for every
      \(A\in\mathcal S\);
\item \(A\mapsto K(y,A)\) is a probability measure on \((S,\mathcal S)\)
      for every \(y\in T\).
\end{enumerate}
Its second input is a measurable set, not a point of \(S\).
Write \(\int f(x)K(y,dx)\) for integration with respect to that
probability measure.
\index{probability kernel}
\end{hbdefinition}

\begin{hbdefinition}{def:p15-regular-conditional-law}{A regular conditional law of a random element}
Let \(X:(\Omega,\mathcal F_0,P)\to(S,\mathcal S)\) be measurable and
\(\mathcal G\subseteq\mathcal F_0\) a sub-sigma-algebra.
A regular conditional law is a probability kernel \(K\) from
\((\Omega,\mathcal G)\) to \((S,\mathcal S)\) with
\[
                      K(\cdot,A)=P(X\in A\mid\mathcal G)
                                  \quad\text{almost everywhere}
                                      \qquad(A\in\mathcal S).
\]
The exceptional set in this equality can initially depend on \(A\).
The source permits its measure-valued condition itself to fail on a null
set. The constructions here specify a probability measure at every
argument by a measurable default on a common exceptional set.

When \(S=\Omega\), \(\mathcal S=\mathcal F_0\) and \(X\) is the identity,
this is a regular conditional probability on the original outcome space.
That special case needs an existence hypothesis on the outcome space itself.
\index{regular conditional distribution}\index{regular conditional probability}
\end{hbdefinition}

\begin{hbtheorem}{thm:p15-kernel-expectation}{proved=annex:p15.kernel-expectation}{Integrating a function against a conditional law}
For a probability kernel \(K\) and measurable \(f:S\to[0,\infty]\),
the function \(y\mapsto\int f\,dK(y,\cdot)\) is measurable.
If \(K\) is a regular conditional law of \(X\) given \(\mathcal G\), then
\[
                         E(f(X)\mid\mathcal G)
                              =\int f(x)K(\cdot,dx).
\]
For real measurable \(f\) with \(E|f(X)|<\infty\), the same formula
holds on the full-measure set where the absolute kernel integral is
finite. Extending the signed integral by zero on its measurable
exceptional set gives an integrable \(\mathcal G\)-measurable version.
\end{hbtheorem}

\section{Constructing the law from countably many conditional probabilities}
\label{sec:p15-conditional-cdf-construction}

\begin{hbtheorem}{thm:p15-real-conditional-law}{proved=annex:p15.real-conditional-law}{A regular law for a real random variable}
For every real random variable \(W\) and every \(\mathcal G\),
a regular conditional law of \(W\) given \(\mathcal G\) exists.
It may be constructed by choosing versions
\[
                          g_q=P(W\leq q\mid\mathcal G)
                                     \qquad(q\in\mathbb Q),
\]
and, on one \(\mathcal G\)-measurable full-measure set, setting
\[
                          F_\omega(t)=\inf_{q\in\mathbb Q,\ q>t}g_q(\omega).
\]
The proof supplies all monotonicity, right-continuity and endpoint-limit
conditions needed for \(F_\omega\) to be a cumulative distribution function.
On the common exceptional set use the cumulative function of \(\delta_0\).
The Stieltjes construction gives the measure, and a generating-class
argument proves its measurability and the conditional identity for every
Borel event.

Any two such kernels agree as probability measures outside one
\(\mathcal G\)-measurable null set, simultaneously for every Borel event.
\end{hbtheorem}

\begin{hbtheorem}{thm:p15-standard-conditional-law}{proved=annex:p15.standard-conditional-law}{Transport through a Borel real coding}
Suppose \((S,\mathcal S)\) admits a measurable isomorphism
\(\phi:S\to E\), where \(E\) is a Borel subset of \(\mathbb R\).
Then every \(S\)-valued random element \(X\) has a regular conditional
law given every sub-sigma-algebra \(\mathcal G\).
In particular this holds for the Borel subsets of complete separable
metric spaces covered by \cref{thm:p15-metric-borel-representation}.

The construction applies the real theorem to \(\phi(X)\), proves
concentration on \(E\) outside one conditioning-null set, and transports
the probability back through \(\phi^{-1}\). A fixed point mass on \(S\)
specifies every remaining argument. Two resulting kernels agree outside
one null set as measures on all of \(\mathcal S\).
This supplies the source's standard Borel case with its image and
concentration requirements explicit.
\end{hbtheorem}

\begin{hbtheorem}{thm:p15-density-conditional-kernel}{proved=annex:p15.density-conditional-kernel}{A regular kernel from a joint density}
For real \(X,Y\) with joint Borel density \(f\geq0\), put
\(m(y)=\int f(x,y)\,dx\) and \(D=\{y:0<m(y)<\infty\}\).
Fix \(x_0\in\mathbb R\). Then
\[
 K(y,A)=
 \begin{cases}
  \displaystyle\frac{\int_Af(x,y)\,dx}{m(y)},&y\in D,\\
  \mathbf1_A(x_0),&y\notin D
 \end{cases}
\]
is a probability kernel, and \(K(Y,\cdot)\) is a regular conditional
law of \(X\) given \(\sigma(Y)\).
The exceptional set of observation values has zero probability under
the law of \(Y\). This one kernel supplies a probability measure for
all target events.
\end{hbtheorem}

\section{A kernel on the observation space}
\label{sec:p15-observation-kernel}

\begin{hbtheorem}{thm:p15-value-conditional-law}{proved=annex:p15.value-conditional-law}{Conditional laws parametrized by the observed value}
Let \(X\) take values in a standard Borel space \((S,\mathcal S)\),
and let \(Y\) take values in a measurable space \((T,\mathcal T)\).
There exists a probability kernel \(K:T\times\mathcal S\to[0,1]\)
such that
\[
                  K(Y,A)=P(X\in A\mid\sigma(Y))
                               \quad\text{almost everywhere}
                                       \qquad(A\in\mathcal S).
\]
If \(\mu_Y\) is the law of \(Y\), then
\[
               P(X\in A,\ Y\in B)=\int_B K(y,A)\,d\mu_Y(y)
                                  \qquad(A\in\mathcal S,\ B\in\mathcal T).
\]
Two such probability kernels agree on every target event outside one
\(\mu_Y\)-null set of observation values.
The source's Theorem~4.1.18, with both variables in the same standard
Borel space, is included. The proof also permits an arbitrary measurable
conditioning space \(T\), using the explicit scalar factorization theorem.
\end{hbtheorem}

\begin{hbtheorem}{thm:p15-kernel-joint-measure}{proved=annex:p15.kernel-joint-measure}{A joint measure and integration with a kernel}
For any probability measure \(\mu\) on \((T,\mathcal T)\) and
probability kernel \(K\) from \(T\) to \(S\), there is a unique
probability measure \(Q\) on
\((T\times S,\mathcal T\otimes\mathcal S)\) with
\[
                              Q(B\times A)=\int_BK(y,A)\,d\mu(y).
\]
For nonnegative product measurable \(h\),
\[
                \int h\,dQ
                      =\int_T\left(\int_Sh(y,x)K(y,dx)\right)d\mu(y),
\]
and the inner integral is a measurable function of \(y\).
For signed integrable \(h\), use its signed inner integral where its
absolute inner integral is finite, and zero elsewhere.

For the conditional kernel above, this \(Q\) is the actual joint law
of \((Y,X)\). If \(E|h(Y,X)|<\infty\), the inner-integral function
therefore gives a version of \(E(h(Y,X)\mid Y)\).
\end{hbtheorem}

\begin{hbtheorem}{thm:p15-conditional-holder}{proved=annex:p15.conditional-holder}{Conditional H\"older through the joint regular law}
Let \(1<p,q<\infty\), \(1/p+1/q=1\), and suppose
\(E|U|^p<\infty\) and \(E|V|^q<\infty\).
Then
\[
 E(|UV|\mid\mathcal G)
    \leq \bigl(E(|U|^p\mid\mathcal G)\bigr)^{1/p}
          \bigl(E(|V|^q\mid\mathcal G)\bigr)^{1/q}
                      \quad\text{almost everywhere}.
\]
The proof uses the regular conditional law of the pair \((U,V)\)
and the already proved ordinary H\"older inequality, as requested in
the source's Exercise~4.1.4.
\end{hbtheorem}

\section{Null observations and the choice of coordinates}
\label{sec:p15-conditioning-coordinate-examples}

\begin{hbdefinition}{def:p15-sphere-coordinate-law}{The specified spherical coordinate model}
Use longitude \(\Theta\in[0,2\pi)\) and latitude
\(\Phi\in(-\pi/2,\pi/2]\), with joint coordinate density
\[
                         \frac{\cos\phi}{4\pi}.
\]
This is the normalized spherical area law in the parametrization
\[
 r(\theta,\phi)
       =(\cos\phi\cos\theta,\ \cos\phi\sin\theta,\ \sin\phi).
\]
For the alternative great-circle coordinates, let
\(\Theta'=\Theta\) on \(\{\Theta<\pi\}\) and
\(\Theta'=\Theta-\pi\) otherwise. Let
\[
 \Psi=
 \begin{cases}
  \Phi,&\Theta<\pi,\\
  \pi-\Phi,&\Theta\geq\pi,\ \Phi\geq0,\\
  -\pi-\Phi,&\Theta\geq\pi,\ \Phi<0 .
 \end{cases}
\]
These coordinates describe the same sphere point through \(r(\Theta',\Psi)\).
Coordinate duplications at the poles and interval boundaries have probability
zero in the specified law.
\end{hbdefinition}

\begin{hbtheorem}{thm:p15-conditioning-coordinate-examples}{proved=annex:p15.conditioning-coordinate-examples}{The sphere calculation and two approaches to the same null event}
In the spherical model, \(\Theta,\Phi\) are independent, with densities
\(1/(2\pi)\) and \(\cos\phi/2\) on their declared intervals.
The alternative \(\Theta',\Psi\) are independent, with densities
\[
                          \frac1\pi,\qquad\frac{|\cos\psi|}{4}
                                          \quad(-\pi<\psi\leq\pi).
\]
Thus a conditional law of the great-circle angle given its longitude
has the latter nonuniform density.

There is also the following explicit null-event example.
For independent uniform \(X,Y\) on \([0,1]\), put \(Z=Y/(1+X)\).
The events \(\{Y=0\}\) and \(\{Z=0\}\) are identical and null.
A version obtained by the joint-density formula gives the uniform
conditional \(X\)-density at \(Y=0\), with mean \(1/2\).
A version obtained from the \((X,Z)\) density at \(Z=0\) gives
\[
                       \frac23(1+x)\mathbf1_{[0,1]}(x),
                       \qquad\text{mean }5/9.
\]
For every \(0<\varepsilon<1/2\), these are also exactly the conditional
densities given \(Y<\varepsilon\) and \(Z<\varepsilon\), respectively.
The two shrinking-event families have the same intersection but different
conditional laws. Both calculations follow their declared observation
variables; no ratio conditional on a null event has been defined.
\end{hbtheorem}

\begin{proofboundary}{def:p15-conditional-laws-next}{Sequential information and martingales}
The scalar conditional calculus and regular kernels are now available for
filtrations, martingales and sequential inference. Their convergence and
stopping theorems require separate proofs and their stated integrability
or boundedness conditions. The existence result here is for the specified
standard Borel target class and is not an assertion for every measurable
outcome space.
\end{proofboundary}
